% =====================================================================
%  PART II -- FOUNDATIONS
%
%  The mathematical spine.  Read in order: the complex is constructed
%  before it is used, the sheaf is laid on it, the operator that measures
%  everything is built, and the final chapter corrects a type error the
%  earlier ones deliberately carry.
% =====================================================================

\chapter{Where the Complex Comes From}\label{ch:ingest}

\section{The problem: $\CC$ may not be assumed}

Every account of sheaves on a complex begins with the complex already
present. For a theory of cognition this is not available: if the base complex
is postulated, then every structural claim the theory makes is a claim about
a hand-drawn picture, and the vertices --- the regions of
\S\ref{sec:hypothesis} --- are exactly as arbitrary as the person drawing
them.

So this chapter comes first. It asks: given input of an arbitrary nature,
what canonically determines a simplicial complex? The answer wanted must have
three properties.

\begin{enumerate}[label=(P\arabic*)]
  \item\label{P1} \textbf{Minimal input.} It should require as little structure
    on the input as possible --- in particular no metric, since there is no
    reason to believe raw input carries one.
  \item\label{P2} \textbf{No free parameters.} A construction with knobs makes
    every downstream invariant a function of the knobs.
  \item\label{P3} \textbf{Stability.} A small perturbation of the input must
    produce a small perturbation of the complex. Without this, absorbing new
    material is not a controlled operation and no invariant computed after
    ingest means anything.
\end{enumerate}

\section{A relation determines a complex}

\begin{intu}\Intu
The weakest possible thing you can say about a body of input is: \emph{these
things occur together, witnessed by that}. A passage of text witnesses the
co-occurrence of the terms in it; an observation witnesses the co-occurrence
of the features it exhibits. That is a relation between two sets --- things on
one side, witnesses on the other --- and nothing more.

It turns out a relation is already enough. Declare a collection of things to
be jointly bound when \emph{some single witness witnesses all of them at
once}. This is Dowker's construction, and everything the base complex is
required to be follows from it rather than being assumed.
\end{intu}

\begin{definition}[Dowker complexes of a relation]\label{def:dowker}
Let $X,Y$ be finite sets and $R\subseteq X\times Y$ a relation, written
$x\rel y$. Define two abstract simplicial complexes:
\[
  \Dow{X}(R)=\bigl\{\,\sigma\subseteq X \ :\ \sigma\neq\emptyset,\
     \exists\,y\in Y \text{ with } x\rel y \text{ for all } x\in\sigma \,\bigr\},
\]
\[
  \Dow{Y}(R)=\bigl\{\,\tau\subseteq Y \ :\ \tau\neq\emptyset,\
     \exists\,x\in X \text{ with } x\rel y \text{ for all } y\in\tau \,\bigr\}.
\]
Call $\Dow{X}(R)$ the \emph{ingest complex} of $R$ and $Y$ its
\emph{witness set}.
\end{definition}

\begin{proposition}[Downward closure is free]\label{prop:closurefree}
\Proved\ $\Dow{X}(R)$ is an abstract simplicial complex: it is downward
closed.
\end{proposition}
\begin{proof}
Let $\sigma\in\Dow{X}(R)$ with witness $y$, and let
$\emptyset\neq\tau\subseteq\sigma$. Every $x\in\tau$ lies in $\sigma$, hence
$x\rel y$. So $y$ witnesses $\tau$ and $\tau\in\Dow{X}(R)$.
\end{proof}

This is worth pausing on. In the standard presentation, downward closure is a
\emph{modelling assumption} --- one declares that every sub-collection of a
jointly bound fact is itself recorded, and the reader is asked to accept it.
Under Definition~\ref{def:dowker} it is a one-line consequence of what a
witness is. The same holds for the other structural feature usually assumed:

\begin{proposition}[$n$-ary binding is free]\label{prop:naryfree}
\Proved\ In $\Dow{X}(R)$, a $k$-simplex records a single $(k{+}1)$-ary joint
binding and is not determined by the pairwise bindings among its vertices.
Explicitly, there exist relations $R$ for which every pair
$\{x_i,x_j\}\subseteq\sigma$ lies in $\Dow{X}(R)$ while $\sigma$ does not.
\end{proposition}
\begin{proof}
Take $X=\{x_1,x_2,x_3\}$, $Y=\{y_{12},y_{13},y_{23}\}$, and let $y_{ij}$ be
related to exactly $x_i$ and $x_j$. Each pair has a witness, so all three
edges are present; no single $y$ is related to all of $x_1,x_2,x_3$, so
$\{x_1,x_2,x_3\}\notin\Dow{X}(R)$. Adding a fourth witness $y_{123}$ related
to all three fills the triangle. The two relations have the same
$1$-skeleton and different complexes.
\end{proof}

\begin{remark}[Why this replaces the usual argument]\label{rem:nary}
The usual justification for simplicial complexes over graphs is that a fact
depending on three hypotheses is one three-way fact rather than three pairwise
ones. That argument is correct but weak: it justifies \emph{higher-order}
structure, and a hypergraph supplies that more cheaply and without downward
closure. What it does not justify is the \emph{simplicial} structure
specifically. Propositions~\ref{prop:closurefree} and~\ref{prop:naryfree}
together do: downward closure is not a price paid for the chain complex, it is
what a witness relation already satisfies. The chain complex --- and with it
everything in Chapters~\ref{ch:sheaf} and~\ref{ch:laplacian} --- is then
available for free.
\end{remark}

\section{Dowker duality}

\begin{theorem}[Dowker, 1952]\label{thm:dowker}
\Proved\ (Cited; not proved here.) For any relation $R\subseteq X\times Y$,
the complexes $\Dow{X}(R)$ and $\Dow{Y}(R)$ are homotopy equivalent. In
particular they have isomorphic homology and cohomology in every
degree.~\cite{Dowker}
\end{theorem}

The content for this theory: \emph{ingest does not privilege a side}. The
complex built on the things and the complex built on the witnesses carry the
same topology, so the invariants of Chapter~\ref{ch:laplacian} are invariants
of the relation, not of the arbitrary decision to call one side ``objects''
and the other ``contexts''. A functorial and stable refinement, which is what
\ref{P3} requires, is in Chowdhury--M\'emoli~\cite{ChowdhuryMemoli}.

\begin{hon}\Hon
This is the point at which formal concept analysis is often proposed, since
it also produces structure canonically from a binary relation via a Galois
connection. It is not used here, and the reason is structural rather than one
of taste: the output of that construction is a \emph{lattice}. A lattice has
no boundary operator, hence no cochain complex, hence no cohomology and no
Laplacian --- so none of Chapters~\ref{ch:sheaf}--\ref{ch:laplacian} can be
run on it. The Dowker complex keeps the one property that made the relational
input attractive (canonical, parameter-free, no metric required) and adds the
homotopy-theoretic structure the rest of the book consumes.
\end{hon}

\section{The alternatives, and their prices}\label{sec:alternatives}

\begin{center}
\begin{tabular}{@{}>{\raggedright\arraybackslash}p{0.24\textwidth}%
                  >{\raggedright\arraybackslash}p{0.20\textwidth}%
                  >{\raggedright\arraybackslash}p{0.13\textwidth}%
                  >{\raggedright\arraybackslash}p{0.33\textwidth}@{}}
\toprule
\textbf{construction} & \textbf{input needed} & \textbf{parameters} &
\textbf{stability} \\
\midrule
Dowker complex & a relation & none & yes~\cite{ChowdhuryMemoli} \\
\addlinespace[2pt]
nerve of a cover & a cover of a space & the cover &
  good cover $\Rightarrow$ $\Nrv(\mathcal U)\simeq X$ \\
\addlinespace[2pt]
Vietoris--Rips / \v Cech & a metric & scale $\varepsilon$ &
  yes~\cite{CEH,ChazalStructure} \\
\addlinespace[2pt]
Mapper & metric, filter, cover, clustering & four & poor \\
\bottomrule
\end{tabular}
\end{center}

The nerve construction is the conceptual archetype --- it is where the word
``region'' comes from, and the nerve theorem is the statement that a good cover
loses no topology --- but it presupposes a space to cover, which ingest does
not have. Vietoris--Rips requires a metric on the input and a scale; it is the
right choice when the input genuinely carries a metric, and the persistence
machinery exists precisely to remove the arbitrariness of $\varepsilon$ by
taking all scales at once. Mapper is named because it is the best-known
data-to-complex pipeline and because it fails \ref{P2} and \ref{P3} badly
enough that its output is not a measurement.

\begin{hon}\Hon
Requirement \ref{P3} is the one that is easy to skip and expensive to lose.
The stability theorems~\cite{CEH,ChazalStructure} are what make ``absorb this
text'' a controlled operation: they bound the change in the computed
invariants by the size of the change in the input. Without such a bound, a
one-word change to a document could in principle alter $b_1$ arbitrarily, and
no downstream quantity would be reproducible. Any future ingest operator
proposed for this theory must come with a stability statement or be rejected.
\end{hon}

\section{The ingest operator, and the running example}

\begin{construction}[Ingest]\label{con:ingest}
\Proposed\ Input arrives as a finite relation $R\subseteq X\times Y$ between
a set $X$ of \emph{things} and a set $Y$ of \emph{witnesses}. Ingest returns
the complex $\CC=\Dow{X}(R)$, truncated to dimension $2$
(\S\ref{sec:dimtwo}). Assigning content to its cells is
Chapter~\ref{ch:sheaf}; joining it to an existing complex is
\S\ref{sec:mv}.
\end{construction}

\begin{example}[The running example $\Kc_4$, derived]\label{ex:K4}
Let $X=\{\mathsf a,\mathsf b,\mathsf c,\mathsf d\}$ and $Y=\{y_1,y_2\}$ with
\[
  R=\bigl\{(\mathsf a,y_1),(\mathsf b,y_1),(\mathsf c,y_1),
           (\mathsf c,y_2),(\mathsf d,y_2)\bigr\}.
\]
Read it concretely: two source passages, $y_1$ and $y_2$; the first exhibits
$\mathsf a,\mathsf b,\mathsf c$ together, the second exhibits
$\mathsf c,\mathsf d$ together. Then
\[
  \Kc_4 \;=\; \Dow{X}(R) \;=\;
  \bigl\{\,\mathsf a,\mathsf b,\mathsf c,\mathsf d,\;
    \mathsf{ab},\mathsf{ac},\mathsf{bc},\mathsf{cd},\;
    \mathsf{abc}\,\bigr\},
\]
a filled triangle $\{\mathsf a,\mathsf b,\mathsf c\}$ with a pendant edge
$\{\mathsf c,\mathsf d\}$. The triangle is filled because a single witness
$y_1$ carries all three at once; had the three terms merely occurred in
pairs across three different passages, the triangle would be empty
(Proposition~\ref{prop:naryfree}). Every instrument built in this book is
applied to this object.
\end{example}

\begin{center}
\begin{tikzpicture}[scale=1.6,
  vtx/.style={circle,fill=black,inner sep=1.8pt},
  lbl/.style={font=\small}]
  \coordinate (A) at (0,1);
  \coordinate (B) at (-1,0);
  \coordinate (C) at (1,0);
  \coordinate (D) at (2.3,-0.6);
  \fill[linkblue,opacity=0.18] (A) -- (B) -- (C) -- cycle;
  \draw[thick] (A)--(B) (B)--(C) (A)--(C) (C)--(D);
  \node[vtx] at (A) {}; \node[vtx] at (B) {};
  \node[vtx] at (C) {}; \node[vtx] at (D) {};
  \node[lbl,above] at (A) {$\mathsf{a}$};
  \node[lbl,left]  at (B) {$\mathsf{b}$};
  \node[lbl,below right] at (C) {$\mathsf{c}$};
  \node[lbl,right] at (D) {$\mathsf{d}$};
  \node[font=\footnotesize,linkblue] at (0,0.35) {filled by $y_1$};
\end{tikzpicture}
\captionof{figure}{The running example $\Kc_4=\Dow{X}(R)$ of
Example~\ref{ex:K4}. Nothing about it is chosen: it is what the relation
$R$ says.}
\label{fig:K4}
\end{center}

\begin{workedbox}\Worked
\emph{Dowker duality, checked by hand.} The witness-side complex is
$\Dow{Y}(R)=\{y_1,y_2,\{y_1,y_2\}\}$: the pair $\{y_1,y_2\}$ is a simplex
because $\mathsf c$ is related to both. So $\Dow{Y}(R)$ is a single edge,
which is contractible. And $\Dow{X}(R)$ is a filled triangle with a whisker,
also contractible. Both are contractible, hence homotopy equivalent, in
agreement with Theorem~\ref{thm:dowker}. \hfill$\square$

Note what the check would have caught: if the triangle $\mathsf{abc}$ had been
left \emph{unfilled} while its three edges were present, $\Dow{X}(R)$ would
have $b_1=1$ whereas $\Dow{Y}(R)$ still has $b_1=0$, contradicting the
theorem --- which is another way of seeing that filling is forced by the
relation and not optional.
\end{workedbox}

\begin{agenda}
\agset{
  \item Downward closure (Proposition~\ref{prop:closurefree}) and non-reducibility
    of $n$-ary binding to pairwise binding (Proposition~\ref{prop:naryfree})
    are theorems about relations, not modelling assumptions.
  \item Ingest is side-independent, by Dowker's theorem.
  \item The running example is derived from a relation rather than drawn.
  \item Stability is stated as a \emph{requirement} on any ingest operator, and
    the alternatives are ranked against it.
}
\agopen{
  \item Ingest satisfies \ref{P1} and \ref{P2}, but the choice of the witness
    set $Y$ is itself a modelling decision and is not canonical. The
    construction has no parameters \emph{given} $R$; producing $R$ from raw
    input is unaddressed.
  \item The stability of $\Dow{X}$ is cited~\cite{ChowdhuryMemoli} for the
    persistent setting. The one-scale statement actually used here --- how
    much $\Kc$ and its invariants move when a bounded number of pairs are
    added to or removed from $R$ --- is not stated and not proved.
  \item Truncation to dimension~$2$ (Construction~\ref{con:ingest}) discards
    cells the relation genuinely produces. Its effect on $b_1$ and on
    $\Hone$ is not analysed anywhere.
  \item Nothing here says how a witness set is \emph{grown}. Ingest builds a
    complex from a relation; it does not say what happens to $Y$ when new
    input arrives with new witnesses.
}
\agsteps{
  \item State and prove a one-scale stability bound: for $R,R'$ with
    $|R\,\triangle\,R'|\le k$, bound the change in $b_0$, $b_1$ and in
    $\lambda_2(\Lap)$. This is the missing form of \ref{P3} and everything
    downstream that claims to be a measurement depends on it.
  \item Analyse $2$-truncation. Either bound the error it induces in the
    invariants actually used, or drop the truncation and pay the
    combinatorial cost knowingly.
  \item Give at least one concrete, fully specified $R$ from a non-toy input,
    and compute $\Kc$, $b_0$, $b_1$ on it. Until this exists, ingest is a
    definition rather than an operator.
  \item Decide whether $Y$ is part of the state. If witnesses persist, they
    belong in the theory; if they are consumed at ingest, say so and record
    that the complex retains no trace of what justified its cells.
}
\end{agenda}

% =====================================================================
\chapter{The Base Complex}\label{ch:complex}

\section{Simplicial complexes}

\begin{definition}[Abstract simplicial complex]\label{def:asc}
Let $V$ be a finite set. An \emph{abstract simplicial complex} $\CC$ on $V$
is a family of nonempty subsets of $V$, called \emph{simplices}, which is
downward closed: if $\sigma\in\CC$ and $\emptyset\neq\tau\subseteq\sigma$
then $\tau\in\CC$. A simplex of cardinality $k+1$ is a \emph{$k$-simplex};
write $\dim\sigma=k$. If $\tau\subseteq\sigma$ we write $\tau\triang\sigma$
and call $\tau$ a \emph{face} of $\sigma$. The set of $k$-simplices is
$\CC_k$, and $\dim\CC=\max_\sigma\dim\sigma$.
\end{definition}

$0$-simplices are vertices --- the \emph{regions} of \S\ref{sec:hypothesis};
$1$-simplices are edges, the places where two regions are compared;
$2$-simplices are filled triangles. By Proposition~\ref{prop:closurefree} the
closure condition is automatic for any complex arising from ingest.

\begin{example}[The one-region space]\label{ex:atom}
\Proved\ Let $\CC=\{\{v\}\}$, a single vertex. Then $\CC_1=\CC_2=\emptyset$,
so with any sheaf whatsoever (Chapter~\ref{ch:sheaf}):
\[
  \Cone=0,\qquad \cob=0,\qquad \Lap=0,\qquad
  \discord(x)=0 \text{ for all } x,\qquad \Hzero=\stalk{v}.
\]
Every content is a thought; nothing can fail; every invariant the book
defines is zero or trivial.
\end{example}

This is the collapse promised in \S\ref{sec:hypothesis}, and it is the precise
sense in which the theory is a theory of the relations. Content lives in the
stalks; \emph{all} structure lives in the maps between them. If meaning
admitted a global chart, Example~\ref{ex:atom} would be the whole theory.

\section{Why dimension two, and why not one}\label{sec:dimtwo}

The book works with complexes of dimension at most~$2$. Both bounds need
justification, and they are of very different kinds.

\paragraph{Not dimension one.} A graph is a $1$-dimensional simplicial
complex, so working on a graph is a strict special case, not a different
choice. What is lost is stated exactly in Theorem~\ref{thm:functoriality} (the
sheaf condition becomes vacuous, so nothing can obstruct) and
Theorem~\ref{thm:hodge} (the three-way decomposition loses a summand, so
repairable misalignment and irreparable holes stop being distinguishable).
These are the two load-bearing results of Part~\ref{part:foundations} and
they are both statements that dimension~$1$ is degenerate.

\paragraph{Not dimension three or more.} This bound is pragmatic and is not
claimed to be principled. Cells of dimension~$3$ would make $\cob^2$
non-vacuous in turn and would give $H^2$ a role. Nothing in this book needs
them, the combinatorial cost of downward closure grows as $2^{k+1}$ per
$k$-simplex, and no phenomenon has been identified that requires them.

\begin{hon}\Hon
\Contested\ The asymmetry above should be uncomfortable. The argument against
dimension~$1$ is a theorem; the argument against dimension~$3$ is that nothing
has needed it yet. If a later construction turns out to require an obstruction
living in $H^2$, the truncation is what will be blocking it, and the
truncation was chosen for convenience. This is recorded here rather than
discovered later.
\end{hon}

\section{On the word ``stratified''}\label{sec:stratword}

Earlier presentations of this material described $\CC$ as \emph{stratified},
meaning that it carried two nested scales --- a coarse complex whose vertices
were subsystems, and inside each, a fine complex whose vertices were concepts.
That description is withdrawn, for two independent reasons.

First, the word was wrong. In mathematics a stratification is a decomposition
into locally closed smooth pieces satisfying frontier and tangency conditions
(Whitney), or the base data for constructible sheaves and the exit-path
category. None of those theorems were being invoked, and none of them apply to
a finite complex with no smooth structure. Under the naming discipline of
\S\ref{sec:naming}, rule~2, the term is retired.

Second, and more seriously, the structure it named was not earned. Two nested
scales were justified by the claim that the coarse vertices were genuinely
separate subsystems. Under the reading of \S\ref{sec:hypothesis} vertices are
charts, not subsystems, and that ground is gone. Rather than re-justify a
two-scale structure on a new basis, this book takes the opposite position:

\begin{quote}\itshape
$\CC$ has one scale. Multi-scale structure, where it exists, is to be
\emph{derived} from the sheaf --- for instance by spectral partition of $\Lap$
--- and not postulated.
\end{quote}

\begin{remark}[Where a genuine stratification does appear]\label{rem:stratmoved}
The word is not lost, only moved. There is a real stratification in this
theory, but it is on the space of \emph{configurations} $(\CC,\Fs)$ rather
than on $\CC$: the strata are indexed by the dimension vector, molding moves
within a stratum, and growth crosses to an adjacent one.
\S\ref{sec:stratconf} develops this. It is a better home for the word because
there the strata do real work --- they are what makes the identity invariant
locally constant.
\end{remark}

\section{The running example, combinatorially}\label{sec:running}

For $\Kc_4$ of Example~\ref{ex:K4}, with vertex order
$(\mathsf a,\mathsf b,\mathsf c,\mathsf d)$, the degrees are
$2,2,3,1$ and the ordinary graph Laplacian of the $1$-skeleton is
\begin{equation}\label{eq:LK4}
  L_{\Kc_4}=
  \begin{pmatrix}
     2 & -1 & -1 &  0\\
    -1 &  2 & -1 &  0\\
    -1 & -1 &  3 & -1\\
     0 &  0 & -1 &  1
  \end{pmatrix}.
\end{equation}

\begin{proposition}[Spectrum of the running example]\label{prop:specK4}
\Proved\ $\det(L_{\Kc_4}-\lambda I)=\lambda(\lambda-1)(\lambda-3)(\lambda-4)$,
so $\operatorname{spec}L_{\Kc_4}=\{0,1,3,4\}$.
\end{proposition}
\begin{proof}
Direct expansion gives
$\lambda^4-8\lambda^3+19\lambda^2-12\lambda=\lambda(\lambda-1)(\lambda-3)(\lambda-4)$.
The eigenvalue $0$ is forced since every row sums to zero, so
$L_{\Kc_4}\mathbf 1=0$; its multiplicity is $1$ because the $1$-skeleton is
connected (Proposition~\ref{prop:identityfatal}).
\end{proof}

The integrality is an accident of this graph and should not be expected in
general. It is why $\Kc_4$ was chosen: every quantity in
Part~\ref{part:foundations} can be carried through in exact arithmetic.

\begin{agenda}
\agset{
  \item $\CC$ is a single-scale finite complex of dimension $\le 2$; the
    two-scale ``stratified'' description is withdrawn with its reason.
  \item The degenerate one-region case is computed, and it is what the
    founding hypothesis denies.
  \item The lower dimension bound is justified by two theorems; the upper one
    is flagged as pragmatic.
}
\agopen{
  \item Multi-scale structure is asserted to be derivable by spectral
    partition. No derivation is given, no criterion says how many parts, and
    nothing shows the derived coarse complex is again a Dowker complex of
    anything.
  \item The dimension-$2$ truncation is unjustified except by convenience
    (recorded as \Contested).
}
\agsteps{
  \item Define the coarse-graining operator explicitly --- a spectral or
    modular partition of $\CC$ inducing a simplicial quotient --- and state
    what it preserves. Until then, ``derived, not postulated'' is a promise.
  \item Check whether the derived coarse complex admits a witness relation
    presenting it; if it does, coarse-graining and ingest are compatible and
    the theory has one notion of region, not two.
  \item Identify one phenomenon requiring $H^2$, or record that none was found
    and close the dimension question.
}
\end{agenda}

% =====================================================================
\chapter{The Sheaf}\label{ch:sheaf}

\begin{intu}\Intu
On top of the shape we lay data. Each region carries a space of things
expressible in that region's own coordinates. Each comparison site carries a
space in which the neighbouring regions' contents are to be compared, and each
region has a translation into it. A cellular sheaf is exactly this. A thought
is then a choice of content in every region that survives every translation.
\end{intu}

\section{Two algebras, not one}\label{sec:twoalgebras}

Before any definition, a distinction that earlier presentations of this
material collapsed and that costs real theorems if collapsed.

A stalk must do two jobs, and they are answerable to different algebra:

\begin{center}
\begin{tabular}{@{}>{\raggedright\arraybackslash}p{0.26\textwidth}%
                  >{\raggedright\arraybackslash}p{0.30\textwidth}%
                  >{\raggedright\arraybackslash}p{0.36\textwidth}@{}}
\toprule
 & \textbf{coefficient algebra $\Ralg$} & \textbf{scalar field $\Kfield$}\\
\midrule
carries & content: what can be expressed in a region & the metric: how far two
  contents are from agreeing\\
\addlinespace[2pt]
determines & what an obstruction \emph{means} & whether $\Lap$, relaxation and
  the Hodge decomposition exist\\
\addlinespace[2pt]
candidates & $\Bool$, $\Rnn$, a lattice, a field & $\R$, $\Cx$, $\Hq$ only
  (Proposition~\ref{prop:frobenius})\\
\bottomrule
\end{tabular}
\end{center}

Choosing $\Ralg=\Kfield$ --- taking content to be a vector over the same field
that carries the metric --- is the choice this book makes, but it is a choice
with a price, and \S\ref{sec:tradeoff} states the price. Making it silently is
the error to avoid.

\section{The coefficient algebra: what content is}\label{sec:coefficient}

The general form of the question is settled in the contextuality
literature~\cite{AbramskyBrandenburger}, where a presheaf assigns to each
context the distributions over its joint outcomes, valued in a commutative
semiring $\Ralg$, and \emph{contextuality is precisely the non-existence of a
global section}. That is the same diagram as this book's, over a different
site, and there the semiring is an explicit parameter with consequences:

\begin{center}
\begin{tabular}{@{}>{\raggedright\arraybackslash}p{0.22\textwidth}%
                  >{\raggedright\arraybackslash}p{0.30\textwidth}%
                  >{\raggedright\arraybackslash}p{0.40\textwidth}@{}}
\toprule
$\Ralg$ & content is & obstruction detected\\
\midrule
$\Bool=(\{0,1\},\vee,\wedge)$ & possibility &
  logical (possibilistic) obstruction\\
\addlinespace[2pt]
$\Rnn$ & probability & probabilistic obstruction\\
\addlinespace[2pt]
$\R$, $\Cx$ (signed) & signed amplitude &
  \emph{weakened or absent}: allowing negative weights admits global sections
  that positivity forbids\\
\addlinespace[2pt]
a lattice & propositions &
  order-theoretic obstruction, via a Tarski Laplacian~\cite{GhristRiess}\\
\bottomrule
\end{tabular}
\end{center}

\begin{hon}\Hon
\Contested\ The third row is the important one and it is stated here in the
weakest form I can defend: passing from a positive semiring to a signed one
can destroy obstructions that positivity created, because the constraint being
violated was positivity and not linearity. The precise statement in the source
literature --- for which models, and in what generality, a signed global
section always exists --- must be checked against~\cite{AbramskyBrandenburger}
and~\cite{ABKLM} before it is quoted as a theorem of this book. It is recorded
now because the \emph{direction} of the effect is what motivates
\S\ref{sec:tradeoff}, and that much is not in doubt.
\end{hon}

\section{The scalar field is forced}

\begin{definition}[Discord axiom]\label{def:discordaxiom}
\Proposed\ The disagreement of a state is required to be a single element of
$\R_{\ge0}$, vanishing exactly when there is no disagreement, and minimisable.
\end{definition}

This is a modelling axiom and it is where the real commitment is made ---
not in the classification below, which is only its consequence.

\begin{proposition}[The field is forced to $\R$, $\Cx$ or $\Hq$]
\label{prop:frobenius}
\Proved\ Suppose stalks are finite-dimensional modules over an associative
division algebra $\Kfield$ which is finite-dimensional over $\R$, equipped
with an involution and a Hermitian form $\inner{\cdot}{\cdot}$ for which
$\inner{x}{x}$ satisfies Definition~\ref{def:discordaxiom}. Then
$\Kfield\in\{\R,\Cx,\Hq\}$.
\end{proposition}
\begin{proof}
The hypotheses make $\Kfield$ a finite-dimensional associative division
algebra over $\R$. By Frobenius' theorem the only such algebras up to
isomorphism are $\R$, $\Cx$ and $\Hq$~\cite{Palais}.
\end{proof}

\begin{hon}\Hon
The proposition is one line and the hypotheses are everything. What has been
shown is not that the field must be one of three, but that
\emph{Definition~\ref{def:discordaxiom} costs exactly that much}: demand a
single non-negative real measure of disagreement and the algebra is
classified. Any future proposal to work over another field is a proposal to
give up Definition~\ref{def:discordaxiom}, and should be stated that way.
\end{hon}

\begin{proposition}[The metric theory fails over finite fields]
\label{prop:finitefield}
\Proved\ Let $\Fq$ have odd characteristic, let $n\ge3$, and let $B$ be a
nondegenerate symmetric bilinear form on $\Fq^{\,n}$. Take $\CC$ to be a
single edge $e=\{u,v\}$ with $\stalk{u}=\stalk{v}=\stalk{e}=\Fq^{\,n}$ and
identity restriction maps. Then there is a state $x\neq 0$ with
$\discord(x)=0$ but $\cob x\neq0$.
\end{proposition}
\begin{proof}
The quadratic form $Q(z)=B(z,z)$ in $n\ge3$ variables over a finite field is
isotropic: it has a nontrivial zero~\cite{Lam}. Pick $z\neq0$ with $Q(z)=0$
and set $x_u=0$, $x_v=z$. Then $(\cob x)_e=z\neq0$, so $\cob x\neq0$, while
$\discord(x)=B(z,z)=0$.
\end{proof}

So over $\Fq$ the equivalence $\discord(x)=0\iff\cob x=0$ ---
Proposition~\ref{prop:Lpsd}, on which every later result depends --- is false.
Cohomology survives over any field; the \emph{metric} theory does not.

\section{The trade-off, stated as a cost}\label{sec:tradeoff}

\begin{center}
\begin{tabular}{@{}>{\raggedright\arraybackslash}p{0.30\textwidth}%
                  >{\raggedright\arraybackslash}p{0.30\textwidth}%
                  >{\raggedright\arraybackslash}p{0.32\textwidth}@{}}
\toprule
 & \textbf{$\Kfield$-vector stalks} & \textbf{lattice stalks}~\cite{GhristRiess}\\
\midrule
Laplacian & $\Lap=\cob^{\adjt}\cob$ & Tarski Laplacian\\
$\Hzero$ & $\ker\Lap$ & fixed points\\
convergence & spectral, rate $\lambda_2$ & monotone, by Tarski\\
three-way Hodge split & \textbf{yes} & no\\
positivity obstructions & \textbf{lost} & preserved\\
\bottomrule
\end{tabular}
\end{center}

This book takes the left column. The reason is Theorem~\ref{thm:hodge}: the
three-way decomposition is what distinguishes repairable misalignment from
irreparable holes, and it is what makes the routing of
\S\ref{sec:degenerate} possible at all. The right column keeps a class of
obstruction that the left column destroys.

\begin{hon}\Hon
This is a real loss and it is placed here rather than in a limits chapter so
that it cannot be forgotten. Cognitive Space Theory is a \emph{metric} theory
of coherence, and it purchases its dynamics by giving up possibilistic
obstruction. Should a phenomenon be found that is purely possibilistic ---
incoherent in the $\Bool$ semiring but coherent over $\Kfield$ --- this theory
will not see it, by construction and not by oversight.
\end{hon}

\begin{hon}\Hon
\Proposed\ Between $\R$ and $\Cx$, this book's recommendation is $\Cx$, and
the reasons are cheap: every result below holds verbatim with transpose
replaced by conjugate transpose; restriction maps become unitary rather than
orthogonal, and $U(d)$ is connected whereas $O(d)$ is not, so the
reflection-parity obstruction to continuously adjusting maps disappears;
phase permits asymmetric relations inside a self-adjoint operator, removing
the need for a second directed structure alongside $\CC$; and the
contextuality literature this theory borrows from is complex-valued. The text
below is written for a general $\Kfield\in\{\R,\Cx\}$ and specialises to $\R$
in worked examples only.
\end{hon}

\section{Which way do the restriction maps point?}\label{sec:direction}

\begin{proposition}[The direction is forced]\label{prop:direction}
\Proved\ Suppose we want a linear operator $\cob$ from vertex data to edge
data whose kernel is exactly the set of assignments $x=(x_v)_{v\in V}$ that
agree on every edge. Then the structure maps must go from faces to cofaces,
$\rest{u\triang e}:\stalk{u}\to\stalk{e}$.
\end{proposition}
\begin{proof}
An assignment gives one vector $x_v\in\stalk{v}$ per vertex and nothing on
edges. For the endpoint values $x_u,x_v$ of $e=\{u,v\}$ to be compared they
must be brought into a common space, and the only space available to $e$ is
$\stalk{e}$. Hence maps $\stalk{u}\to\stalk{e}$, $\stalk{v}\to\stalk{e}$ are
required, and agreement on $e$ reads
$\rest{u\triang e}(x_u)=\rest{v\triang e}(x_v)$.

If instead the maps ran downward, $\rho_{e\to u}:\stalk{e}\to\stalk{u}$, then
from vertex data no map into $\stalk{e}$ exists, so no operator
$\Czero\to\Cone$ is constructible from the structure maps and the condition
``$x$ agrees on $e$'' is not expressible. Downward maps assemble instead into
a boundary operator $\partial:\Cone\to\Czero$ whose kernel is edge data
summing to zero at vertices --- a flow condition, not an agreement condition.
That is a cellular \emph{cosheaf}, and its invariants are homology, not the
cohomology used here.
\end{proof}

\section{The definition}

\begin{definition}[Cellular sheaf]\label{def:sheaf}
A \emph{cellular sheaf} $\Fs$ of finite-dimensional $\Kfield$-inner-product
spaces on $\CC$ assigns
\begin{enumerate}[label=(\roman*)]
  \item to each $\sigma\in\CC$ a finite-dimensional inner-product space
        $\stalk{\sigma}$ over $\Kfield$, the \emph{stalk};
  \item to each face relation $\sigma\triang\tau$ a linear
        \emph{restriction map}
        $\rest{\sigma\triang\tau}:\stalk{\sigma}\to\stalk{\tau}$,
\end{enumerate}
subject to $\rest{\sigma\triang\sigma}=\Id$ and, for every chain
$\sigma\triang\tau\triang\upsilon$,
\begin{equation}\label{eq:functoriality}
  \rest{\tau\triang\upsilon}\circ\rest{\sigma\triang\tau}
  \;=\;\rest{\sigma\triang\upsilon}.
\end{equation}
Equivalently, $\Fs$ is a functor from the face poset of $\CC$ to
$\Vect_{\Kfield}$.
\end{definition}

\begin{remark}[The inner product is where confidence lives]
\label{rem:metricmeaning}
``Inner-product space'' is not a technical convenience. The inner product on
$\stalk{\sigma}$ is a free modelling parameter: rescaling it reweights that
cell's contribution to the disagreement, and two sheaves with the same stalks
and the same restriction maps but different stalk metrics have
\emph{different Laplacians and different coherent states}. Precision or
confidence, wherever the theory needs it, is carried here and nowhere else.
\end{remark}

\begin{definition}[Cochains]\label{def:cochains}
For $k\ge0$ set $\Ck{k}=\bigoplus_{\sigma\in\CC_k}\stalk{\sigma}$, the
\emph{$k$-cochains}, with the direct-sum inner product
$\inner{x}{y}=\sum_\sigma\inner{x_\sigma}{y_\sigma}$.
\end{definition}

\section{What the $2$-cells do: functoriality is $\cob^2=0$}

\begin{theorem}[Functoriality on a triangle is exactly $\cob^1\cob^0=0$]
\label{thm:functoriality}
\Proved\ Let $\tau=\{u,v,w\}\in\CC_2$ with edges $uv,uw,vw$. With the
coboundary maps of Definition~\ref{def:coboundary} and
\eqref{eq:coboundary2}, the following are equivalent:
\begin{enumerate}[label=(\alph*)]
  \item $(\cob^1\cob^0x)_\tau=0$ for every $x\in\Czero$;
  \item the composite $\stalk{p}\to\stalk{\tau}$ is independent of the
    intermediate edge, for each vertex $p\in\tau$ --- that is,
    \eqref{eq:functoriality} holds on $\tau$.
\end{enumerate}
Consequently, if $\dim\CC\le1$ the condition is vacuous and \emph{every}
assignment of stalks and linear maps is a sheaf.
\end{theorem}
\begin{proof}
Write $\rho_{p\triang e}$ for vertex-to-edge and $\rho_{e\triang\tau}$ for
edge-to-triangle maps. Expanding,
\[
  (\cob^1\cob^0x)_\tau
   = \rho_{vw\triang\tau}(\cob x)_{vw}
   - \rho_{uw\triang\tau}(\cob x)_{uw}
   + \rho_{uv\triang\tau}(\cob x)_{uv},
\]
and substituting $(\cob x)_{pq}=\rho_{q\triang pq}x_q-\rho_{p\triang pq}x_p$
and collecting the coefficient of each of $x_u,x_v,x_w$ gives
\begin{align*}
  (\cob^1\cob^0x)_\tau
  &= \bigl[\rho_{uw\triang\tau}\rho_{u\triang uw}
          -\rho_{uv\triang\tau}\rho_{u\triang uv}\bigr]x_u\\
  &\quad+\bigl[\rho_{uv\triang\tau}\rho_{v\triang uv}
          -\rho_{vw\triang\tau}\rho_{v\triang vw}\bigr]x_v\\
  &\quad+\bigl[\rho_{vw\triang\tau}\rho_{w\triang vw}
          -\rho_{uw\triang\tau}\rho_{w\triang uw}\bigr]x_w .
\end{align*}
Since $x_u,x_v,x_w$ range independently over their stalks, the expression
vanishes identically if and only if each bracket vanishes, which is exactly
(b). If $\dim\CC\le1$ there are no chains $\sigma\triang\tau\triang\upsilon$
of length two, so \eqref{eq:functoriality} imposes no condition.
\end{proof}

This is the central structural fact of Part~\ref{part:foundations}, and it is
what the choice of a $2$-complex buys. Its consequences:

\begin{corollary}[On a graph nothing can obstruct]\label{cor:graphvacuous}
\Proved\ If $\dim\CC\le1$ the set of sheaves on $\CC$ with fixed stalk
dimensions is the \emph{full} affine space
$\prod_{v\triang e}\operatorname{Hom}(\stalk v,\stalk e)$. Any continuous
adjustment of restriction maps is therefore unobstructed: it can be carried
out in any direction, from any starting point, without leaving the space of
sheaves.
\end{corollary}
\begin{proof}
Immediate from Theorem~\ref{thm:functoriality}: with no constraint, every
tuple of linear maps is a sheaf.
\end{proof}

\begin{remark}[Why this is the whole argument for growth]
\label{rem:growthderivable}
On a $2$-complex the sheaves form the subvariety of
$\prod\operatorname{Hom}$ cut out by the quadratic equations of
Theorem~\ref{thm:functoriality}(b). Adjusting restriction maps --- molding ---
is motion \emph{within} that variety. When the data demand a configuration
outside it, no adjustment reaches the demand, and the only remaining move is
to change $\CC$ itself. That is the sense in which growth becomes
\emph{derivable} rather than triggered by a hand-written rule, and by
Corollary~\ref{cor:graphvacuous} it is unavailable on a graph. The clause
``growth from an obstruction'' in the thesis of \S\ref{sec:isnot} rests
entirely on this.
\end{remark}

\begin{hon}\Hon
\Proposed\ It is tempting to quantify: with all stalks of dimension $d$, the
ambient space has $(2E+3T)d^2$ parameters and the constraints supply at most
$3Td^2$ equations, suggesting a subvariety of codimension $3Td^2$. The
equations are \emph{not} obviously independent, the variety is not obviously
smooth, and no such count is asserted here. Establishing the true codimension
--- and whether the variety is connected --- is exactly what would turn
Remark~\ref{rem:growthderivable} from a mechanism into a theorem with a
stated obstruction class.
\end{hon}

\section{Flatness, holonomy, and what a filled triangle is}
\label{sec:flatness}

\begin{proposition}[Invertible sheaves are local systems]\label{prop:localsystem}
\ProvedCond\ Suppose every restriction map is invertible. Then $\Fs$ is a
local system on $\CC$: on each simplex the composite maps define parallel
transport, $\Hzero$ is the space of parallel sections, and on a connected
$\CC$ the sheaf is determined up to isomorphism by the holonomy
representation $\pi_1(|\CC|)\to GL(\stalk{v_0},\Kfield)$ at a basepoint.
\end{proposition}
\begin{proof}[Proof sketch]
Invertibility makes the functor of Definition~\ref{def:sheaf} factor through
the fundamental groupoid of $|\CC|$; \eqref{eq:functoriality} is the statement
that transport is independent of the path within a simplex. The identification
of $\Hzero$ with parallel sections is then
Proposition~\ref{prop:Lpsd}. Full details are standard and are not reproduced;
see~\cite{Curry}.
\end{proof}

\begin{intu}\Intu
Read together with Theorem~\ref{thm:functoriality}: transport around the
boundary of a \emph{filled} triangle is required to be trivial. A filled
triangle is a place where the connection is forced to be flat. An
\emph{unfilled} cycle carries holonomy, and that holonomy is a curvature the
system can measure and cannot remove by re-choosing frames. So the three
structural commitments --- higher cells, non-identity maps, obstruction-driven
growth --- are one geometric statement: \emph{fill where flatness is asserted;
measure where it is not; grow where the measurement will not go away.}
\end{intu}

\section{The potency ladder, and why the identity is fatal}

Restriction maps admit three levels of expressiveness:

\begin{center}
\begin{tabular}{@{}>{\raggedright\arraybackslash}p{0.10\textwidth}%
                  >{\raggedright\arraybackslash}p{0.26\textwidth}%
                  >{\raggedright\arraybackslash}p{0.56\textwidth}@{}}
\toprule
level & maps & what it can express\\
\midrule
0 & identity & ``the same thing, verbatim'' --- no knowledge in the wiring\\
1 & unitary / orthogonal & ``the same content in a different frame'' ---
  reversible, lossless\\
2 & general linear & ``agreeing only on a subspace'' --- lossy, and the loss
  is itself a measurement of the connection\\
\bottomrule
\end{tabular}
\end{center}

Level~0 is not a simplification but a deletion, and
Proposition~\ref{prop:identityfatal} makes that precise once the Laplacian is
available.

\begin{agenda}
\agset{
  \item Content algebra and metric field are separated, and the metric field
    is classified by Proposition~\ref{prop:frobenius} given
    Definition~\ref{def:discordaxiom}.
  \item Finite fields are excluded by an explicit counterexample
    (Proposition~\ref{prop:finitefield}).
  \item The cost of the metric choice --- loss of possibilistic obstruction ---
    is stated where the choice is made.
  \item The direction of restriction maps is forced
    (Proposition~\ref{prop:direction}).
  \item Theorem~\ref{thm:functoriality}: functoriality on $2$-cells is exactly
    $\cob^2=0$, and it is vacuous on graphs.
}
\agopen{
  \item The recommendation $\Kfield=\Cx$ is argued but not adopted: the book's
    worked examples remain real, and no result below has been re-derived in
    the complex case even though all are expected to hold verbatim.
  \item The signed-semiring claim in \S\ref{sec:coefficient} is
    \Contested\ and must be verified against its sources before use.
  \item Remark~\ref{rem:growthderivable} describes a mechanism, not a theorem.
    No obstruction class is exhibited; nothing yet \emph{computes} that a
    required configuration lies outside the sheaf variety.
  \item Nothing here says how stalk dimensions are chosen, or what fixes the
    stalk metrics of Remark~\ref{rem:metricmeaning}.
  \item Proposition~\ref{prop:localsystem} is proved only in sketch, and only
    under invertibility, which level~2 of the potency ladder abandons.
}
\agsteps{
  \item Settle $\Kfield$. If $\Cx$, restate Chapters~\ref{ch:sheaf}
    and~\ref{ch:laplacian} with conjugate transposes and re-run the worked
    example; if $\R$, record why the arguments for $\Cx$ were declined.
  \item Verify the signed-semiring statement against
    \cite{AbramskyBrandenburger,ABKLM} and either promote it to \Proved\ with
    a precise hypothesis or withdraw it.
  \item Determine the codimension and connectedness of the sheaf variety on a
    $2$-complex. This is the single highest-value open problem in
    Part~\ref{part:foundations}: it converts growth from a mechanism into a
    derivation.
  \item Give a criterion for stalk dimension and stalk metric, or state that
    both are inputs.
  \item Write out Proposition~\ref{prop:localsystem} in full, and state what
    survives when restriction maps are merely injective, or general.
}
\end{agenda}

% =====================================================================
\chapter{Coboundary, Laplacian, and the Three Channels}\label{ch:laplacian}

\section{The coboundary}

\begin{definition}[Coboundary]\label{def:coboundary}
Fix an arbitrary orientation of each simplex. The degree-zero coboundary
$\cob^0:\Czero\to\Cone$ acts edgewise on $e=(u,v)$ by
\begin{equation}\label{eq:coboundary}
  (\cob^0 x)_e\;=\;\rest{v\triang e}(x_v)\;-\;\rest{u\triang e}(x_u)
  \ \in\ \stalk{e},
\end{equation}
and the degree-one coboundary $\cob^1:\Cone\to\Ctwo$ acts on
$\tau=(u,v,w)$ by
\begin{equation}\label{eq:coboundary2}
  (\cob^1 y)_\tau\;=\;\rest{vw\triang\tau}(y_{vw})
   -\rest{uw\triang\tau}(y_{uw})+\rest{uv\triang\tau}(y_{uv}).
\end{equation}
A $0$-cochain with $\cob^0x=0$ is a \emph{global section}. We write $\cob$
for $\cob^0$ where no ambiguity arises.
\end{definition}

\begin{lemma}[Orientation-independence]\label{lem:orientation}
\Proved\ Reversing the orientation of a simplex replaces the corresponding
component of $\cob x$ by its negative. Consequently $\ker\cob$,
$\norm{\cob x}^2$ and $\cob^{\adjt}\cob$ are independent of the orientations
chosen.
\end{lemma}
\begin{proof}
Reversal exchanges the roles of the endpoints in \eqref{eq:coboundary},
negating that entry. Writing $\cob'=D\cob$ with $D$ blockwise diagonal with
entries $\pm1$, we get $\cob'^{\adjt}\cob'=\cob^{\adjt}D^{\adjt}D\cob
=\cob^{\adjt}\cob$ since $D^{\adjt}D=I$; vanishing and squared norms are
likewise unaffected.
\end{proof}

\begin{remark}
This licenses ``fix an arbitrary orientation'', and it is worth checking
rather than assuming, because the analogous statement fails for the
\emph{signed} per-edge quantity $(\cob x)_e$, which is defined only up to
sign. Any diagnostic reporting $(\cob x)_e$ rather than
$\norm{(\cob x)_e}^2$ is reporting an orientation-dependent quantity and is
not well defined.
\end{remark}

\section{The Laplacian and discord}

\begin{definition}[Sheaf Laplacian and discord]\label{def:laplacian}
The degree-zero sheaf Laplacian is
$\Lap=\cob^{\adjt}\cob:\Czero\to\Czero$, and the \emph{discord} of a state
$x$ is
\begin{equation}\label{eq:omega}
  \discord(x)=\inner{x}{\Lap x}=\norm{\cob x}^2
  =\sum_{e=\{u,v\}}
   \norm{\rest{v\triang e}(x_v)-\rest{u\triang e}(x_u)}^2 .
\end{equation}
\end{definition}

\begin{proposition}[Basic properties]\label{prop:Lpsd}
\Proved\ $\Lap$ is self-adjoint and positive semidefinite, and
\[
  \discord(x)=0\iff\cob x=0\iff x\in\ker\Lap,
\]
with $\ker\cob=\Hzero$.
\end{proposition}
\begin{proof}
Self-adjointness: $(\cob^{\adjt}\cob)^{\adjt}=\cob^{\adjt}\cob$. Positive
semidefiniteness: $\inner{x}{\Lap x}=\norm{\cob x}^2\ge0$. If $\Lap x=0$ then
$\norm{\cob x}^2=0$, so $\cob x=0$ because $\norm{\cdot}$ is a norm --- this
is the step that fails over a finite field
(Proposition~\ref{prop:finitefield}); conversely $\cob x=0$ gives $\Lap x=0$.
Finally the cochain complex begins in degree $0$, so
$\Hzero=\ker\cob^0/\im\cob^{-1}=\ker\cob$.
\end{proof}

\begin{intu}\Intu
$\discord$ is one number: the total squared disagreement across all
comparison sites. It vanishes exactly when there is a single coherent thought.
The achievement of the construction is that ``a coherent thought'' has become
a precise object --- a vector in the kernel of a self-adjoint operator ---
rather than a mood.
\end{intu}

\section{The three channels}

\begin{theorem}[Hodge decomposition]\label{thm:hodge}
\Proved\ Let $\Fs$ be a cellular sheaf on a complex of dimension $\le2$ over
$\Kfield\in\{\R,\Cx\}$, so that $\cob^1\cob^0=0$
(Theorem~\ref{thm:functoriality}). Then
\begin{equation}\label{eq:hodge}
  \Cone \;=\; \underbrace{\im\cob^0}_{\text{gradient}}
        \;\oplus\; \underbrace{\Harm^1}_{\text{harmonic}}
        \;\oplus\; \underbrace{\im(\cob^1)^{\adjt}}_{\text{curl}},
  \qquad
  \Harm^1:=\ker\cob^1\cap\ker(\cob^0)^{\adjt},
\end{equation}
the sum is orthogonal, and $\Harm^1\cong\Hone$. If $\dim\CC\le1$ then
$\cob^1=0$ and the third summand is zero.
\end{theorem}
\begin{proof}
Orthogonal complements give $\Cone=\im\cob^0\oplus\ker(\cob^0)^{\adjt}$ and
$\Cone=\ker\cob^1\oplus\im(\cob^1)^{\adjt}$. Since $\cob^1\cob^0=0$ we have
$\im\cob^0\subseteq\ker\cob^1$, so intersecting the first decomposition with
$\ker\cob^1$ gives
$\ker\cob^1=\im\cob^0\oplus(\ker\cob^1\cap\ker(\cob^0)^{\adjt})
=\im\cob^0\oplus\Harm^1$. Substituting into the second decomposition yields
\eqref{eq:hodge}, and orthogonality is inherited from the two complements.
Finally $\Hone=\ker\cob^1/\im\cob^0\cong\Harm^1$ by the first identity. If
$\dim\CC\le1$ then $\Ctwo=0$, so $\cob^1=0$ and $\im(\cob^1)^{\adjt}=0$.
\end{proof}

\begin{intu}\Intu
This is the taxonomy of \S\ref{ch:picture} made exact. A disagreement pattern
$y\in\Cone$ decomposes uniquely and orthogonally into three parts, and each
has a different repair:
\begin{itemize}
  \item \textbf{gradient}: $y=\cob x$ for some state. The disagreement is an
    artefact of the current state and is removed by \emph{relaxing the state}.
  \item \textbf{curl}: in $\im(\cob^1)^{\adjt}$. Circulation that the filled
    triangles see. Removable by \emph{adjusting the restriction maps}.
  \item \textbf{harmonic}: in $\Harm^1\cong\Hone$. Neither. It is a
    topological invariant of the complex, and the only way to remove it is to
    \emph{change the complex}.
\end{itemize}
On a graph the curl channel is empty, so misalignment and holes are
indistinguishable and there is nothing for molding to discharge. That is the
first row of the table in \S\ref{sec:degenerate}, now proved.
\end{intu}

\section{The degenerate corner, proved}

\begin{proposition}[Linear feedforward networks are path sheaves]
\label{prop:nnsheaf}
\Proved\ Let $W_i\in\Kfield^{\,n_{i+1}\times n_i}$ for $i=0,\dots,L-1$. Let
$\CC$ be the path complex on vertices $v_0,\dots,v_L$ with edges
$e_i=\{v_i,v_{i+1}\}$ and no $2$-cells, and define a sheaf by
$\stalk{v_i}=\Kfield^{\,n_i}$, $\stalk{e_i}=\Kfield^{\,n_{i+1}}$,
\[
  \rest{v_i\triang e_i}=W_i,\qquad \rest{v_{i+1}\triang e_i}=\Id .
\]
Then for $x\in\Czero$,
\[
  \discord(x)=\sum_{i=0}^{L-1}\norm{x_{i+1}-W_ix_i}^2,
\]
so $x\in\Hzero$ if and only if $x_{i+1}=W_ix_i$ for all $i$ --- that is, if
and only if $x$ is the trace of a forward pass. Moreover the map
$\Hzero\to\Kfield^{\,n_0}$, $x\mapsto x_0$, is a linear isomorphism, so
$\dim\Hzero=n_0$.
\end{proposition}
\begin{proof}
Substituting the stated maps into \eqref{eq:coboundary} gives
$(\cob x)_{e_i}=x_{i+1}-W_ix_i$, and \eqref{eq:omega} gives the displayed
discord. By Proposition~\ref{prop:Lpsd}, $x\in\Hzero$ iff every term
vanishes, which is the recursion. The evaluation map is linear; it is
injective because the recursion determines $x_1,\dots,x_L$ from $x_0$, and
surjective because any $x_0$ extends by the recursion. The sheaf axioms are
vacuous by Theorem~\ref{thm:functoriality} since $\dim\CC=1$.
\end{proof}

So the set of valid activation traces of a linear network is literally
$\Hzero$ of a path-shaped cellular sheaf, and the forward pass is the unique
global section extending a given input.

\begin{hon}\Hon
The nonlinearity is not absorbed. A pointwise $\sigma$ makes the restriction
maps nonlinear and leaves $\Vect_{\Kfield}$ entirely. Two routes exist ---
sheaves valued in lattices with a Tarski Laplacian~\cite{GhristRiess}, or
enlarging the stalks to graphs of the nonlinearity --- and neither is
developed in this book. Proposition~\ref{prop:nnsheaf} is exact for the linear
case and is not claimed beyond it.
\end{hon}

\begin{proposition}[Identity restriction maps empty the theory]
\label{prop:identityfatal}\label{prop:b0kernel}%
% prop:b0kernel is the name this result carried in the earlier draft; the
% alias is kept so that later parts, not yet rewritten, still resolve.
\Proved\ Suppose every stalk is $\Kd$ and every restriction map is the
identity. Then
\[
  \Lap=L_{\CC}\otimes I_d,
\]
where $L_\CC$ is the ordinary graph Laplacian of the $1$-skeleton. Hence
$\operatorname{spec}\Lap$ is $\operatorname{spec}L_\CC$ with each eigenvalue
repeated $d$ times; $\dim\ker\Lap=d\cdot b_0$; and on a connected $\CC$,
\[
  \Hzero=\{x: x_u=x_v \ \forall u,v\},\qquad
  e^{-t\Lap}x\ \xrightarrow[t\to\infty]{}\ \text{the constant assignment at }
  \tfrac{1}{|V|}\textstyle\sum_v x_v .
\]
\end{proposition}
\begin{proof}
With identity restrictions \eqref{eq:omega} reads
$\discord(x)=\sum_{\{u,v\}\in E}\norm{x_u-x_v}^2$, the Dirichlet form of
$L_\CC$ applied coordinatewise in each of the $d$ components, so
$\Lap=L_\CC\otimes I_d$. The spectrum of a Kronecker product with the identity
is that of the factor with multiplicity $d$. The kernel of $L_\CC$ is spanned
by indicator vectors of connected components, of which there are $b_0$;
tensoring with $\Kd$ multiplies dimension by $d$. On a connected complex the
kernel is spanned by $\mathbf 1$, so $\Hzero$ is the diagonal, and
$e^{-t\Lap}\to$ the orthogonal projection onto $\ker\Lap$, which is the mean.
\end{proof}

\begin{hon}\Hon
This is the sharpest single argument in the book, because it is a theorem and
not a preference. With identity maps: the only coherent thought is exact
consensus; relaxation is plain averaging; and diffusion provably drives every
region to the global mean regardless of where it started. That last statement
is the phenomenon known as oversmoothing, and the published repair is exactly
to make the restriction maps non-trivial~\cite{BodnarSheaf}. Everything a
sheaf can express that a graph cannot --- that two regions agree \emph{after}
a change of frame, or agree only on a subspace --- lives in non-identity
restriction maps. Freezing them at the identity is not a cheap approximation;
it is a proof that the sheaf is decoration.
\end{hon}

\section{The running example, computed}

\begin{workedbox}\Worked
Take $\Kc_4$ (Example~\ref{ex:K4}) over $\Kfield=\R$ with identity
restrictions and $d=2$, and the state
\[
  x_{\mathsf a}=\begin{pmatrix}1\\0\end{pmatrix},\quad
  x_{\mathsf b}=\begin{pmatrix}1\\0\end{pmatrix},\quad
  x_{\mathsf c}=\begin{pmatrix}0\\1\end{pmatrix},\quad
  x_{\mathsf d}=\begin{pmatrix}0\\1\end{pmatrix}.
\]

\emph{Discord, edge by edge}, using $\discord_e=\norm{x_u-x_v}^2$:
\[
\begin{array}{lcl}
  \{\mathsf a,\mathsf b\}:&\ \norm{(0,0)}^2 &= 0,\\
  \{\mathsf b,\mathsf c\}:&\ \norm{(1,-1)}^2 &= 2,\\
  \{\mathsf a,\mathsf c\}:&\ \norm{(1,-1)}^2 &= 2,\\
  \{\mathsf c,\mathsf d\}:&\ \norm{(0,0)}^2 &= 0,
\end{array}
\qquad\Longrightarrow\qquad \discord(x)=4 .
\]

\emph{Localisation.} The two guilty edges tie at $\discord_e=2$, so the
$\arg\max$ is not unique: a diagnostic returning a single ``worst edge'' here
returns an artefact of tie-breaking order, not a measurement. The honest
return is the tied set
$\{\{\mathsf b,\mathsf c\},\{\mathsf a,\mathsf c\}\}$, which reads back
correctly as ``$\mathsf c$ is the region at odds with the rest''.

\emph{Kernel projection.} The $1$-skeleton is connected, so by
Proposition~\ref{prop:identityfatal} the kernel is the diagonal and the
projection is the constant assignment at the mean
$\bar x=\tfrac14\bigl((1,0)+(1,0)+(0,1)+(0,1)\bigr)=(\tfrac12,\tfrac12)$.
Hence
\[
  \norm{x_0}^2=4\cdot\bigl(\tfrac14+\tfrac14\bigr)=2,\qquad
  x_\perp=x-x_0=\bigl(\pm\tfrac12,\mp\tfrac12\bigr)_v,\qquad
  \norm{x_\perp}^2=4\cdot\tfrac12=2,
\]
and Pythagoras checks: $\norm{x}^2=2+2=4$. \hfill$\square$
\end{workedbox}

\begin{exercise}
Recompute $\discord(x)$ after replacing one restriction map by a
non-identity: let $\rest{\mathsf c\triang e}$ be rotation by $90^\circ$ on the
two edges at $\mathsf c$ inside the triangle, and the identity elsewhere. Does
$\discord$ rise or fall --- and does the resulting assignment satisfy
\eqref{eq:functoriality} on the filled triangle
$\{\mathsf a,\mathsf b,\mathsf c\}$? The second half is the important half,
and Theorem~\ref{thm:functoriality} says exactly what to check.
\end{exercise}

\begin{agenda}
\agset{
  \item $\Lap$ is self-adjoint PSD and $\Hzero=\ker\Lap$
    (Proposition~\ref{prop:Lpsd}): a coherent thought is a matrix kernel.
  \item Theorem~\ref{thm:hodge}: three orthogonal channels, with the third
    existing only above dimension~$1$.
  \item Proposition~\ref{prop:nnsheaf}: linear feedforward networks are
    exactly $\Hzero$ of a path sheaf.
  \item Proposition~\ref{prop:identityfatal}: identity maps force consensus
    and averaging, and the published oversmoothing diagnosis is this statement.
}
\agopen{
  \item The routing law --- relax on gradient, mold on curl, grow on harmonic
    --- is stated as intuition. No dynamical system is defined, no convergence
    is proved, and nothing shows the three repairs do not interfere.
  \item $\dim\Hzero\ge1$ is \emph{not} automatic once restriction maps are
    non-identity. A sheaf with $\Hzero=0$ admits no coherent thought but the
    zero one, and no condition is given that prevents molding from arriving
    there.
  \item Neither the harmonic component nor the curl component is computed
    anywhere, including on the running example, because with identity maps on
    $\Kc_4$ both are zero. The book has therefore never exhibited a nonzero
    element of the channel its architecture depends on.
  \item Nothing here treats the normalised Laplacian, and the convergence rate
    $\lambda_2$ is named without being bounded for any family of interest.
}
\agsteps{
  \item Compute a nonzero harmonic class. Build the smallest sheaf on
    $\Kc_4$ with non-identity maps for which $\Harm^1\neq0$ and print the
    class. Until this exists, the three-channel decomposition is a theorem
    with no worked instance, which is the weakest state a load-bearing result
    can be in.
  \item State the routing law as a dynamical system and prove it decreases
    $\discord$; identify its fixed points.
  \item Find a condition on the restriction maps guaranteeing
    $\dim\Hzero\ge1$, and check whether molding preserves it. If it does not,
    molding can destroy the possibility of coherent thought and must be
    constrained.
  \item Bound $\lambda_2(\Lap)$ for sheaves arising from ingest, so that
    ``relaxation converges'' has a rate.
}
\end{agenda}

% =====================================================================
\chapter{Configuration, Time, and Gluing}\label{ch:config}

The four preceding chapters carry a deliberate type error, and this chapter
corrects it.

\section{The error in ``the current state''}

It is natural to define a cognitive space as a triple $(\CC,\Fs,s)$ with $s$
a distinguished $0$-cochain called the current state. The word
\emph{current} presupposes a time parameter that the triple does not contain,
and adding one does not fix the problem, because the space $s$ lives in is
itself time-dependent: $\Czero=\bigoplus_v\stalk{v}$ changes whenever $\CC$
grows or a stalk dimension changes. A trajectory of $s$ is therefore
\emph{not} a curve in a fixed vector space, and writing $s(t)\in\Czero$
asserts something false.

\section{Configuration space, and the stratification}\label{sec:stratconf}

\begin{construction}[The configuration space]\label{con:config}
\Proposed\ Let $\Conf$ be the space whose points are pairs $(\CC,\Fs)$ ---
a finite complex of dimension $\le2$ together with a cellular sheaf on it.
Over each point sits the vector space $\Czero$, and these assemble into a
bundle $p:\mathcal{E}\to\Conf$. A \emph{history} is a path $\gamma$ in
$\Conf$ together with a section of $\gamma^*\mathcal{E}$.
\end{construction}

\begin{definition}[The dimension vector]\label{def:Kinv}
For a configuration $(\CC,\Fs)$ put
$\Kinv(\CC,\Fs)=\bigl(\dim\stalk{\sigma}\bigr)_{\sigma\in\CC}$.
\end{definition}

\begin{construction}[Stratification by $\Kinv$]\label{con:strat}
\Proposed\ Decompose $\Conf$ by the value of $\Kinv$. Within a stratum the
complex and all stalk dimensions are fixed and only the restriction maps vary,
so a stratum is the sheaf variety of \S\ref{sec:flatness} --- the subvariety
of $\prod\operatorname{Hom}$ cut out by Theorem~\ref{thm:functoriality}(b).
Then:
\begin{itemize}
  \item \textbf{molding} is motion within a stratum, and $\Kinv$ is constant
    along it \emph{by the type of the motion}, not by a separate argument;
  \item \textbf{growth} is a transition to an adjacent stratum, and it changes
    $\Kinv$ by an amount determined by which transition was taken.
\end{itemize}
\end{construction}

This is where the word retired in \S\ref{sec:stratword} belongs. The strata
here do work: they are why a dimension-vector invariant is \emph{locally
constant}, which is the standard reason invariants are discrete.

\begin{hon}\Hon
\Proposed\ Construction~\ref{con:strat} is a proposal with the right shape and
no proofs. $\Conf$ has been given no topology; ``adjacent'' has no definition;
the strata are not known to be smooth, connected, or even nonempty for a given
$\Kinv$; and the sheaf variety's structure is exactly the open problem
recorded in Chapter~\ref{ch:sheaf}. What is claimed is only that this is the
right \emph{place} for the two-phase separation to live, and that placing it
here would convert three separately-argued claims into one structural fact.
\end{hon}

\section{Continuous time}

\begin{definition}[Sheaf heat flow]\label{def:heat}
Within a fixed stratum, the relaxation of the state is the flow
\begin{equation}\label{eq:heat}
  \frac{ds}{dt}=-\Lap s,\qquad s(t)=e^{-t\Lap}s(0).
\end{equation}
\end{definition}

\begin{proposition}[Relaxation converges to a thought]\label{prop:heatconv}
\Proved\ $e^{-t\Lap}s(0)\to P_{\ker\Lap}s(0)$ as $t\to\infty$, where
$P_{\ker\Lap}$ is the orthogonal projection onto $\Hzero$, and the rate is
governed by the smallest nonzero eigenvalue $\lambda_2$ of $\Lap$:
$\norm{e^{-t\Lap}s-P_{\ker\Lap}s}\le e^{-\lambda_2 t}\norm{s-P_{\ker\Lap}s}$.
\end{proposition}
\begin{proof}
$\Lap$ is self-adjoint PSD (Proposition~\ref{prop:Lpsd}), so it has an
orthonormal eigenbasis with eigenvalues $0=\lambda_1\le\lambda_2\le\dots$.
In that basis $e^{-t\Lap}$ acts by $e^{-t\lambda_i}$, which is $1$ on
$\ker\Lap$ and tends to $0$ elsewhere, with the slowest decay $e^{-\lambda_2t}$
on $(\ker\Lap)^\perp$.
\end{proof}

The discrete update $s_{t+1}=s_t-\eta\Lap s_t$ is the explicit Euler scheme
for \eqref{eq:heat}; the continuous form is the primary object and the
discrete one an approximation to it, not the other way round.

\begin{hon}\Hon
Two things are being suppressed and should not be.

First, if $\Hzero=0$ then \eqref{eq:heat} converges to $0$: the only coherent
state is silence and relaxation destroys the state. This is the open item
O2 of Chapter~\ref{ch:laplacian} appearing as a dynamical failure rather than
an algebraic curiosity.

Second, taking the state to be a \emph{point} of $\Czero$ rather than a
\emph{distribution} on it is a modelling choice with an alternative: a density
evolving by the Fokker--Planck equation associated with \eqref{eq:heat} has as
its stationary measure the Gaussian with precision $\Lap$, a Gaussian Markov
random field. That version carries uncertainty natively. It is not developed
here, and the choice is recorded as a choice.
\end{hon}

\section{Gluing, and a candidate for the growth address}\label{sec:mv}

Chapter~\ref{ch:ingest} produces a complex from input; this section joins it
to what is already present, and asks what the join creates.

\begin{construction}[Anchor and glue]\label{con:glue}
\Proposed\ Let $\CC$ be the present complex and $\CC'$ the ingest complex of
new input. An \emph{anchoring} is a span $\CC\hookleftarrow A\hookrightarrow\CC'$
of subcomplex inclusions together with an isomorphism of the two restricted
sheaves on $A$. The glued complex is the pushout $\CC\cup_A\CC'$, carrying the
sheaf determined by the two restrictions and the isomorphism.
\end{construction}

The anchoring is the only free choice in the pipeline of
\S\ref{sec:reading}, and it is where interpretation lives: two anchorings of
the same input into the same complex are two readings of the same text.

\begin{proposition}[Mayer--Vietoris for cellular sheaves]\label{prop:mv}
\Proved\ Let $\CC=A\cup B$ with $A,B$ subcomplexes and let $\Fs$ be a sheaf
on $\CC$. Then the sequence of cochain complexes
\[
  0\to C^\bullet(\CC;\Fs)\to
  C^\bullet(A;\Fs|_A)\oplus C^\bullet(B;\Fs|_B)\to
  C^\bullet(A\cap B;\Fs|_{A\cap B})\to 0
\]
given by restriction and difference is exact, and hence induces a long exact
sequence
\[
  \cdots\to H^0(\CC;\Fs)\to H^0(A)\oplus H^0(B)
  \xrightarrow{\ -\ } H^0(A\cap B)
  \xrightarrow{\ \partial\ } H^1(\CC;\Fs)\to\cdots
\]
\end{proposition}
\begin{proof}
Cochains are direct sums over cells, and the cells of $A\cup B$ are the union
of those of $A$ and $B$, with the cells of $A\cap B$ counted in both. Exactness
is then inclusion--exclusion cellwise: a cochain on $\CC$ is precisely a pair
of cochains on $A$ and $B$ agreeing on $A\cap B$ (exactness in the middle and
injectivity on the left), and every cochain on $A\cap B$ extends to $A$ by
zero elsewhere (surjectivity on the right). The maps commute with $\cob$
because restriction does. The long exact sequence is the standard consequence.
\end{proof}

\begin{intu}\Intu
Read the connecting map $\partial$. An element of $H^0(A\cap B)$ is a coherent
thought \emph{on the overlap}. If it is not the difference of coherent
thoughts on the two sides, then $\partial$ sends it to a nonzero class in
$H^1(\CC)$: the act of gluing has created a hole that neither piece
had. That is what it means for reading to change you --- \emph{the new
obstruction is a property of the join, not of either part}. A text you already
agreed with entirely contributes nothing to the image of $\partial$; a text
you cannot anchor at all has empty $A\cap B$ and also contributes nothing.
\end{intu}

\begin{hon}\Hon
\Proposed\ This is offered as a \emph{candidate} for the growth address and
not as a solution. What recommends it is that its type is right: the theory
needs a $1$-dimensional obstruction that is measured from data rather than
manufactured from the current state, and a class in the image of $\partial$ is
by construction not a coboundary. What is missing is substantial. The map
$\partial$ depends on the anchoring of Construction~\ref{con:glue}, and
nothing says how it varies with that choice; no procedure turns a class in
$H^1$ into a specific cell to add; and the relation between this class and the
harmonic component of Theorem~\ref{thm:hodge} --- which is where the routing
law expects to find the address --- is stated nowhere.
\end{hon}

\section{The definition, restated}

\begin{definition}[Cognitive space]\label{def:M}
A \emph{cognitive space} over $\Kfield$ is a pair
$\M=(\CC,\Fs)$ consisting of a finite abstract simplicial complex $\CC$ of
dimension at most $2$ (Definition~\ref{def:asc}) and a cellular sheaf $\Fs$ of
finite-dimensional $\Kfield$-inner-product spaces on $\CC$
(Definition~\ref{def:sheaf}). A \emph{state} of $\M$ is an element
$s\in\Czero$; it is \emph{coherent} when $\cob s=0$. A \emph{history} is a
path in the configuration space $\Conf$ (Construction~\ref{con:config})
together with a section of the pulled-back cochain bundle along it.
\end{definition}

The state is not part of the space. That is the whole content of the
correction: $\M$ is the structure, $s$ is a point of a vector space attached
to it, and the two move on different footings --- $s$ by
Proposition~\ref{prop:heatconv} within a fixed configuration, $\M$ by molding
within a stratum and growth across strata.

\begin{agenda}
\agset{
  \item The state is separated from the space, and the type error in
    ``current state'' is identified and corrected.
  \item Relaxation has a continuous form with proved convergence and rate
    (Proposition~\ref{prop:heatconv}).
  \item Mayer--Vietoris for cellular sheaves is proved
    (Proposition~\ref{prop:mv}) and supplies a candidate of the correct type
    for the growth address.
}
\agopen{
  \item $\Conf$ has no topology, so Construction~\ref{con:config} is a
    bookkeeping device rather than a space and ``path in $\Conf$'' is not yet
    meaningful.
  \item Construction~\ref{con:strat} is entirely unproved: no smoothness, no
    connectivity, no definition of stratum adjacency.
  \item The dependence of $\partial$ on the anchoring is unexamined, and no
    procedure converts a class in $H^1$ into a cell to add.
  \item The relation between the Mayer--Vietoris class and the harmonic
    component of Theorem~\ref{thm:hodge} is not established, although the
    routing law requires them to be the same object.
  \item The Fokker--Planck alternative to a point-valued state is recorded and
    not developed.
}
\agsteps{
  \item Topologise $\Conf$. The natural candidate within a stratum is the
    subspace topology from $\prod\operatorname{Hom}$; between strata, define
    adjacency by elementary growth moves and check that the resulting relation
    is a partial order.
  \item Prove or disprove that the image of $\partial$ has nonzero harmonic
    part. This is the load-bearing question: if the classes created by gluing
    are harmonic, the routing law and the ingest pipeline are the same
    mechanism, and if they are not, one of the two is wrong.
  \item Determine how $\partial$ varies with the anchoring, and identify which
    downstream quantities are anchoring-independent. Anything that is not
    invariant is a property of the reader, not of the text, and must be
    labelled as such.
  \item Define the growth move: given a harmonic class, which cell is added.
    Until this exists, ``growth is derivable'' remains a claim about where the
    trigger comes from and says nothing about what the trigger does.
}
\end{agenda}
